\documentclass[11pt,a4paper]{article}
\usepackage[T1]{fontenc}
\usepackage{lmodern,amsmath,amssymb,booktabs,microtype}
\usepackage[margin=24mm]{geometry}
\usepackage[colorlinks=true,urlcolor=blue,linkcolor=blue]{hyperref}
\newcommand{\EE}{\mathbb E}
\newcommand{\tr}{\operatorname{tr}}
\newcommand{\PP}{\mathsf P_h}
\newcommand{\MA}{\mathsf A}
\newcommand{\BB}{\mathsf B}
\newcommand{\FF}{\mathsf F}
\newcommand{\Jh}{\mathcal J_h}
\newcommand{\W}{\mathsf W}
\title{Two components of the finite-time memory fluctuation\\A conserved four-field correction and a mixed spectral check}
\author{Alan Day \and Chengheng\thanks{AI research assistant (OpenAI).}}
\date{Local research note --- 7 November 2026}
\begin{document}\maketitle
\begin{abstract}
The four-field term in the Gaussian Dyson expansion of the ordinary
memory force has a strong leading energy regression. Independent
auxiliary copies, three-color chord moments and the standard
$q$-Gaussian conditional expectation prove this in annealed $L^2$,
uniformly on compact real-time intervals. Exact energy pinching
therefore removes its collective contribution. The actual memory
field reduces to two bounded components: a charged-transfer defect
and a neutral covariance-filter residual. We identify the first two
neutral Taylor coefficients and their mixed spectral measures with
the charged coefficients. A nonzero second-order cross is a necessary
calibration for the still-unknown joint finite-time kernel.
\end{abstract}

\section{Model, state and previously established input}
Use standard CAR, $\hbar=1$, string real time, and the overlapping
complex $p$-body Gaussian model
\begin{align}
 H_N&=\sum_{I,J}J_{IJ}c_I^\dagger c_J, &
 \EE J_{IJ}J_{KL}&=w_N\delta_{IL}\delta_{JK},\\
 w_N&=4^p\frac{J_s^2}{\lambda_N}\left(\frac{p!}{N^p}\right)^2,
 &\lambda_N&=p^2/N.
\end{align}
All charges share the same original $J$. Take even $N$, fixed $J_s>0$,
and first $\lambda_N\to\lambda>0$. The state is the normalized full
$\beta=\mu=0$ annealed trace, with $b_m=2^{-N}\binom Nm$.
All $L^r$ norms include this trace and the disorder average.
At charge $m$ write $h=H_m$ and
\begin{equation}
 C_m=w_N\binom mp\binom{N-m+p}p,\quad
 M=\sup_m C_m\le (J_s^2/\lambda_N)e^{2p/N}.
\end{equation}
Let $P_e$ be the exact spectral projections of $h$ and
$\PP X=\sum_eP_eXP_e$, including whole degenerate blocks.
The preceding stages \cite{Uniform,Return} prove uniform second
moments and $\sup_{|s|\le S}\|\FF(s)-X_0(s)\|_2\to0$ for the
actual memory field and its ordinary-propagation term. Here we
further reduce $X_0$, keeping all common-disorder cross terms.

\newpage
\section{Exact components and the new reduction}
For $m>0$ define
\begin{align}
 V\psi&=m^{-1/2}\sum_i|i\rangle c_i\psi,&
 K&=1_N\otimes H_{m-1},&Q&=1-VV^\dagger,\\
 R&=QKV,&D&=QKQ,&\alpha_m&=(1-p/m)_+,\\
 \Lambda_m(s)&=R^\dagger e^{isK}R,&
 M_m(s)&=R^\dagger e^{isD}R,&\epsilon_m&=1-\alpha_m^2.
\end{align}
The fields are
\begin{equation}
 X_{0,m}=(1-\PP)m e^{-ihs/2}\Lambda_m(s)e^{-ihs/2},\qquad
 \FF_m=(1-\PP)m e^{-ihs/2}M_m(s)e^{-ihs/2}.
\end{equation}
Inactive sectors $m<p$ have zero Hamiltonian and force. At $m=0$
all collective fields are set to zero.
For an independent copy $u$ of $h$, let $\Phi_m(X)=\EE_u[uXu]$ and
\begin{equation}
 \Lambda_{{\rm OU},m}(s)
 =\epsilon_m\EE_u[u e^{is(\alpha_mh+\sqrt{\epsilon_m}u)}u].
\end{equation}
The same $u$ occurs in its three positions. Define, exactly at finite $N$,
\begin{align}
 \MA_m(s)&=(1-\PP)m e^{-ihs/2}
          [\Lambda_m(s)-\Lambda_{{\rm OU},m}(s)]e^{-ihs/2},\\
 \BB_m(s)&=(1-\PP)m\epsilon_m e^{-ihs/2}
          \Phi_m(e^{i\alpha_mhs})e^{-ihs/2}.
\end{align}
Both use the original $h$; auxiliary averaging does not replace
the physical charged force by an independent probe.

The prior even Dyson expansion is
\begin{align}
 \Lambda_{\rm OU}(s)&=\epsilon\Phi(e^{i\alpha hs})
       -\epsilon^2\mathcal I_s^{(\alpha)}(h)+\mathcal E_s,\\
 \mathcal I_s^{(\alpha)}(h)&=s^2\int_{0\le y\le x\le1}
 \EE_u[uU_{s(1-x)}uU_{s(x-y)}uU_{sy}u]\,dx\,dy,
 \quad U_t=e^{i\alpha ht},\\
 \|\mathcal I_s^{(\alpha)}\|_\infty&\le\tfrac32 C_m^2s^2,
 &\|\mathcal E_s\|_2&\le
 \epsilon^3|s|^4\sqrt{10395}\,C_m^3/24.
\end{align}
All three fourth-moment Wick pairings occur in $\mathcal I$.
We prove
\begin{equation}
 \boxed{\quad\sup_{|s|\le S}\|\FF(s)-\MA(s)-\BB(s)\|_2\to0.\quad}
 \label{eq:reduction}
\end{equation}
Each of $\MA,\BB$ is uniformly $L^2$ bounded by the weighted-transfer
and neutral-filter bounds of \cite{Uniform}. Thus their full
four-term covariance sum replaces the actual covariance, uniformly
for compact memory times and all real external times.

\newpage
\section{Why four auxiliary fields have strong energy regression}
Let $h,u,v$ now denote three orthogonal $q$-Gaussian variables with
variance $C=J_s^2/\lambda$ and $q=e^{-4\lambda}$. Write $E_h$ and
$E_{hv}$ for their trace-preserving conditional expectations.
The standard second-quantization construction \cite{BKS} gives
\begin{equation}
 E_{hv}W(h,u)=E_hW(h,u),\qquad W(h,u)\in W^*(h,u).
 \label{eq:square}
\end{equation}
Indeed, the one-particle projection onto $\operatorname{span}(h,v)$
annihilates every $u$ tensor and retains the $h$ tensors. Applied
factorwise to Wick tensors it is the same projection as onto $h$.
For two polynomials $W_1,W_2$ this proves
\begin{equation}
 \tau_q[W_1(h,u)^\dagger W_2(h,v)]
 =\tau_q[(E_hW_1)^\dagger(E_hW_2)].
 \label{eq:replica}
\end{equation}
This identity concerns two whole blocks. It does not assert
classical independence of arbitrary interleaved colored words.

For the finite model let
\begin{equation}
 W_N(h,u)=u f_1(h)u f_2(h)u f_3(h)u,\qquad
 T_N(h)=\EE_u W_N(h,u),
\end{equation}
where the $f_i$ are fixed polynomials. Independence of auxiliary
couplings $u,v$ gives the exact replica identity
\begin{equation}
 \EE_h\tau T_N^\dagger T_N
 =\EE_{h,u,v}\tau W_N(h,u)^\dagger W_N(h,v).
\end{equation}
The three-color joint moments converge to those of $h,u,v$ above.
To see the color extension, perform Gaussian Wick expansion first:
pairings are allowed only between the same color, but every allowed
chord has the same independent $p$-set sums and CAR trace as before.
The fixed-word overlap rule \cite{BNR,Tall} gives $q$ per crossing,
with the stated variance normalization. No extra covariance is
introduced by color; finitely many pairings suffice at fixed degree.
This argument holds for the full normalized trace. The established
central-window trace bridge also gives its fixed-sector version
uniformly for $m=N/2+xp$ with bounded $x$ \cite{Central}.

Let $t(x)$ be the ordinary polynomial defined by $t(h)=E_hW(h,u)$.
Mixed moments with $t(H_N)$ and \eqref{eq:replica} give
\begin{equation}
 \boxed{\quad\|T_N(H_N)-t(H_N)\|_2\longrightarrow0.\quad}
 \label{eq:polynomial}
\end{equation}
It follows by expanding the squared norm: its three limiting terms
are $\|t\|_2^2$, $-2\|t\|_2^2$, and $\|t\|_2^2$.
An annealed mean alone would not establish this cancellation.

\newpage
\section{Full real time, charge weights and exact pinching}
First set $\alpha=1$ and define
\begin{equation}
 T_{abc,N}(h)=\EE_u[u h^a u h^b u h^c u],\qquad
 \mathcal I_s^{(1)}=
 \sum_{n\ge0}\frac{s^2(is)^n}{(n+2)!}
             \sum_{a+b+c=n}T_{abc,N}(h).
 \label{eq:series}
\end{equation}
Gaussian moments and Schatten H\"older with $n+4$ factors give
\begin{equation}
 \|T_{abc,N}\|_2\le\sqrt{(2n+7)!!}\,M^{(n+4)/2}.
\end{equation}
There are $\binom{n+2}{2}$ triples. The degree-$n$ norm is therefore
bounded, for $|s|\le S$, by
\begin{equation}
 \frac{S^{n+2}}{2n!}\sqrt{(2n+7)!!}\,M^{(n+4)/2}.
 \label{eq:tail}
\end{equation}
The sum converges for every fixed $S$, uniformly in $N$.
Thus fixed-degree joint moments are sufficient: first truncate,
then take $N\to\infty$, and finally remove the degree cutoff.

Define $\iota_s(h)$ by replacing the conditional average in
$\mathcal I_s^{(1)}$ with $E_h$ in the $q$-Gaussian algebra.
BKS's $C^*$ version \cite{BKS} makes it a continuous function on
$[-E_*,E_*]$, where $E_*=2\sqrt{C/(1-q)}$. It obeys
\begin{equation}
 \|\iota_s\|_\infty\le\tfrac12s^2 E_*^4.
\end{equation}
Extend it to all real energies by $\widetilde\iota_s(x)=\iota_s(\chi(x))$,
where $\chi$ clips to this interval. For each fixed polynomial
truncation, Hamiltonian moment convergence and uniform integrability
compare its evaluation at $H_N$ with this bounded extension.
The polynomial limit \eqref{eq:polynomial}, the tail \eqref{eq:tail}
and a finite time mesh consequently prove
\begin{equation}
 \sup_{|s|\le S}\|\mathcal I_s^{(1)}(H_N)
                    -\widetilde\iota_s(H_N)\|_2\to0.
 \label{eq:time}
\end{equation}
Time uniformity also follows directly from the normally convergent
series; the bounded target is uniformly Lipschitz on compact $s$.

On $N/4\le m\le3N/4$, $1-\alpha_m\le4p/N$.
Telescoping the three unitaries, using their integral derivatives
and Schatten-$10$ H\"older, bounds
\begin{equation}
 \|\mathcal I_s^{(\alpha_m)}-\mathcal I_s^{(1)}\|_{2,m}
 \le\tfrac12|1-\alpha_m|\,|s|^3\sqrt{9!!}\,M^{5/2}.
\end{equation}
Here $m\epsilon_m^2\le16\lambda_N$. Outside that broad central
interval, $m\epsilon_m^2\le N$ and the uniform operator bound on
$\mathcal I$ combines with the binomial weight $2e^{-N/8}$.
Since the bounded comparison commutes with $h$, \eqref{eq:time} proves
\begin{equation}
 \boxed{\quad\sup_{|s|\le S}
 \|(1-\PP)m\epsilon_m^2e^{-ihs/2}
                 \mathcal I_s^{(\alpha_m)}e^{-ihs/2}\|_2\to0.\quad}
\end{equation}
Likewise $m\mathcal E_s\to0$: centrally
$m\epsilon_m^3\le128\lambda_N^2/p$, and the exterior tail absorbs $N^2$.
Together with the previous return theorem this proves \eqref{eq:reduction}.

\newpage
\section{An explicit expression for the conserved background}
Let $W_n$ be monic variance-$C$ $q$-Hermite polynomials,
\begin{equation}
 xW_n=W_{n+1}+C[n]_qW_{n-1},\qquad [n]_q=(1-q^n)/(1-q).
\end{equation}
Take $v=\alpha h+\beta u$, $\alpha^2+\beta^2=1$ and
$\beta^2=\epsilon$. The left/right Wick tensor recursions of
\cite{BKS}, followed by projection to pure $h$ tensors, give
\begin{equation}
 E_h[uW_n(v)u]=Cq^n\alpha^nW_n(h)
 +\epsilon C^2[n]_q[n-1]_q\alpha^{n-2}W_{n-2}(h).
 \label{eq:OUWick}
\end{equation}
The second term is zero for $n<2$. Only two structures survive:
the external $u$'s pair with each other, or each contracts with
an internal $v$. This keeps the same-copy effect in full.

Set $P_rW_n=r^nW_n$, $D_qW_n=[n]_qW_{n-1}$ and
$\mathcal N_qW_n=nW_n$. Then, for $e_s(x)=e^{isx}$,
\begin{equation}
 \Lambda^q_{\rm OU}(s)=\epsilon CP_{q\alpha}e_s(h)
                +\epsilon^2C^2P_\alpha D_q^2e_s(h).
\end{equation}
This extends from polynomials since the $q$-Gaussian fields are
bounded and $D_q^2$ is bounded on their $L^2$ space.
Its norm is at most $1/[C(1-q)]$ in the normalized Hermite basis.

Expanding along $\alpha=\sqrt{1-\epsilon}$ and comparing with the
frozen-$\alpha h$ Dyson expansion gives
\begin{equation}
 \boxed{\quad
 \iota_s=\frac C2P_q[(\mathcal N_q-x\partial_x)e_s]
                   -C^2D_q^2e_s.\quad}
 \label{eq:iota}
\end{equation}
The $x\partial_x$ term comes from $e_s(\alpha h)$ and cannot be
omitted. $\mathcal N_q$ is the energy-Hermite number operator,
not the coupling-Gaussian chaos number at finite $N$.
Its action on $e_s$ is justified by the bounded Jacobi operator's
factorial coefficient tail.

For a direct polynomial implementation write
$x^n=\sum_jd_{n,j}W_j(x)$, with
\begin{equation}
 d_{0,0}=1,\qquad d_{n+1,j}=d_{n,j-1}+C[j+1]_qd_{n,j+1}.
\end{equation}
If $S_k=\sum_{a+b+c=k}E_h[u h^a u h^b u h^c u]$, then
\begin{equation}
 [S_k]_{W_j}=\tfrac C2(k+2-j)q^j d_{k+2,j}
            +C^2[j+2]_q[j+1]_q d_{k+2,j+2}.
\end{equation}
In particular,
\begin{equation}
 \iota_0=\iota'_0=0,\quad \iota''_0=(2+q)C^2I,\quad
 \iota'''_0=iC^2(1+4q+3q^2+q^3)h.
\end{equation}
These are the leading conserved coefficients. Finite-$N$ $\mathcal I_s$
need not commute with $h$; its dynamic part vanishes only in the
specified double-scaling limit.

\newpage
\section{The first two neutral fluctuation coefficients}
Write $\BB(s)=isB_1-s^2B_2/2+O(s^3)$ at finite $N$ and
$\Jh X=\{h,X\}/2$. Direct differentiation gives
\begin{align}
 B_1&=m\epsilon\alpha(1-\PP)\Phi(h),\\
 B_2&=m\epsilon(1-\PP)[\alpha^2\Phi(h^2)-\alpha\{h,\Phi(h)\}].
\end{align}
For exact covariance eigenvalues $\kappa_r=C_m\mu(r)$, define
\begin{equation}
 Z=h^2-C_mI,\quad h=\sum_a h_a,\quad
 \EE\tr_m h_a^2=C_m\omega_a,\quad
 \bar\kappa=\sum_a\omega_a\kappa_a,\quad B_Z=\Phi(h)-\bar\kappa h.
\end{equation}
The same coefficients can be written exactly as
\begin{align}
 B_1&=m\epsilon\alpha(1-\PP)B_Z,\\
 B_2&=m\epsilon(1-\PP)
       [\alpha^2(\Phi-\kappa_{2p})Z-2\alpha\Jh B_Z].
 \label{eq:B2exact}
\end{align}
Before pinching the difference of the two $B_2/(m\epsilon)$
expressions is the explicitly conserved polynomial
\begin{equation}
 (\alpha^2\kappa_{2p}-2\alpha\bar\kappa)h^2
       +\alpha^2C_m(C_m-\kappa_{2p})I.
\end{equation}

In a central compact window $m=N/2+xp$, the exact product for
$\mu$ gives, for fixed $k,d$,
\begin{equation}
 p[\mu(kp-d)-\mu(kp)]\to4\lambda q^kd,\qquad
 0\le p[\mu(kp-d)-\mu(kp)]\le12\lambda_Nd.
 \label{eq:slope}
\end{equation}
For example, differentiate the logarithm of its two products:
the reciprocal sum is $4p/N+O(p^2/N^2)$ at fixed deficit,
while $\mu(kp)\to q^k$. Input factorial tails and the prior
sixth output-deficit moment \cite{Central,Raw} make this an
$L^2$ multiplier limit, not merely a pointwise rank estimate.
The output squared tail is bounded by a constant times $R^{-4}$
after a cutoff at deficit $R$.

Let $D_2Z=\sum_r(2p-r)P_rZ$, $D_{\rm in}=\sum_a(p-a)h_a$,
$\bar u=\sum_a(p-a)\omega_a$, and
$L_{\rm in}=\W_2\{h,D_{\rm in}-\bar u h\}$.
The exact input/output decomposition and the prior strong bridges are
\begin{align}
 D_2Z&=S_N+L_{\rm in}+2\bar uZ,\qquad
 S_N=\sum_{a,b,r}(a+b-r)P_r\W_2(h_ah_b),\\
 pB_Z&=F/2+o_{L^2}(1),&
 \Jh F&=4\lambda CqL_{\rm in}+o_{L^2}(1),\\
 S_N&=(O_-+O_+)/2+o_{L^2}(1).
\end{align}
Here $O_\pm=\W_2\sum_i[H,c_i]^\dagger[H,c_i]$ and its reverse
ordering, respectively; $F$ is the previously defined linear
rank-filtered field. With $Y_\pm=O_\pm-cZ$,
$c=2\lambda(1-q)/(1+q)$, substitution into \eqref{eq:B2exact} gives
\begin{equation}
 \boxed{\quad B_1\to F,\qquad
 B_2\to d(Y_-+Y_+)+b\Jh F,
 \quad d=4\lambda Cq^2,\quad b=2(q-1).\quad}
 \label{eq:Blimit}
\end{equation}
Exact pinching can be removed from the limit because these prior
Mehler spectral measures have no mass on the energy diagonal.

\newpage
\section{A nonzero mixed spectral calibration}
Use the established two-energy measures $\mu_r$ with Hamiltonian
variance $C$ and chord parameter $q$, normalized to mass one.
Put $e=(E+E')/2$. The needed prior data \cite{Central,Mixed} are
\begin{align}
 d\nu_{(Y_-,Y_+)}&=
 \begin{pmatrix}w&\eta\\\eta&w\end{pmatrix}d\mu_{q^2},
 &d\nu_F&=V_Fd\mu_q,\\
 w&=4C^2\lambda+\frac{16C^2\lambda^2q}{1+q},
 &\eta&=-4C^2\lambda q+\frac{16C^2\lambda^2q}{1+q},\\
 V_F&=64J_s^6q^2,&\nu_{F,Y_\pm}&=0.
\end{align}
The first two charged coefficients $\MA(s)=A_0+isA_1+O(s^2)$
have the previously identified limits
\begin{equation}
 A_0\to Y_-,\qquad A_1\to L-F=R+a\Jh Y_-,\qquad
 a=-\frac{(1-q)^2}{1+q^2},
\end{equation}
where $\nu_{R,F}=\nu_{R,Y_\pm}=0$. Indeed $\BB(0)=0$, the
four-field term starts at $s^2$, and the actual/ordinary memory
fields have identical zeroth and first derivatives at zero.

Equation \eqref{eq:Blimit} yields
\begin{align}
 d\nu_{B_1,B_2}&=b e V_F\,d\mu_q,\\
 d\nu_{B_2}&=2d^2(w+\eta)d\mu_{q^2}+b^2e^2V_Fd\mu_q,\\
 d\nu_{A_0,B_2}&=d(w+\eta)d\mu_{q^2},&
 d\nu_{A_1,B_2}&=ad e(w+\eta)d\mu_{q^2},\\
 \nu_{A_0,B_1}&=\nu_{A_1,B_1}=0.
\end{align}
In particular $w+\eta>0$ for $0<q<1$, so the second-order mixed
measure is nonzero. Its value at zero external time is
$d(w+\eta)>0$. This is a coefficient-level result: any joint
time kernel admitting termwise matching must reproduce it.
We have not exchanged differentiation with an unproved
finite-time large-$N$ limit.

The coefficient measures extend to the full trace without assuming
an arbitrary extension of every named limiting field. For
$K_j=2\cdot10^6(\lambda_Nj^2+\lambda_N^2j^4)$, the broad-central
rank bounds and Gaussian fourth moments give directly
\begin{align}
 \|B_1\|_{2,m}&\le24\lambda_N C_m^{3/2}\sqrt{K_1},\\
 \|B_2\|_{2,m}&\le24\lambda_N C_m^2
                       [\sqrt{2K_2}+2\sqrt{3K_1}].
\end{align}
One uses $\|B_Z\|_2\le C_m^{3/2}(12p/N)\sqrt{K_1}$,
$\|\Jh B_Z\|_2\le\sqrt3 C_m^2(12p/N)\sqrt{K_1}$, and
$\|D_2Z\|_2\le\sqrt2 C_m\sqrt{K_2}$.
All active sectors have the coarser bounds
\begin{equation}
 \|B_1\|_{2,m}\le2pC_m^{3/2},\qquad
 \|B_2\|_{2,m}\le6\sqrt3 pC_m^2.
\end{equation}
Central annuli contribute $O(e^{-2\lambda_-L^2})$, and the far
tail contributes $O(p^2e^{-N/8})$. The prior full-trace bounds on
$Y,L$, together with spectral Cauchy--Schwarz, handle the mixed
terms with $A_0,A_1$. The central coefficient limits therefore give
the displayed full-trace spectral measures.

\newpage
\section{Verification, scope and the next calculation}
Three new scientific outputs accompany this stage.
Exact rational $q$-Fock calculations are compared with an independent
colored-pairing enumeration: 2520 conditional trace equalities,
7200 three-color replica equalities, 66 same-copy OU Wick identities,
and 42 four-field Taylor comparisons. Parameters are
$q=1/4,1/2,3/4$, $C=1,3/2$, and $\alpha=3/5,\beta=4/5$.
The tests include every triple $a+b+c\le6$ and all three
four-field Wick contractions.

Complete CAR coordinate sets in four systems give 24 real-time
integrals, with time quadrature orders 8 and 12 and both signs of
time. The maximum refinement difference is below $2\cdot10^{-15}$.
An interior $N=5,p=2,m=3$ sector has a nonzero finite-$N$ pinched
residual, confirming that the asymptotic regression is not an exact
conserved identity. Initial test systems with $p=1$ or an edge
sector failed to supply this witness; the failed assertion and
source are retained. No scientific output was overwritten.

A third script checks 12 complete-coordinate covariance/coefficient
identities without floating energy pinching. It retains 360 exact
rational slope inequalities and 80/120-digit comparisons for
$\lambda=1/4,1/2,1$, $p=16,32,64,128$, $x=-2,0,2$, $k=1,2$,
and deficits $0,1,2,3,6$. Relative finite-size errors are retained.
These computations check finite identities and conventions; the
uniform limits come from the proofs above. Quadrature refinement
is not interval certification.

\paragraph{Remaining work.}
The finite-time field now has a two-component reduction with the
full covariance $C_{AA}+C_{AB}+C_{BA}+C_{BB}$. Its joint limiting
kernel has not been evaluated. The next calculation must retain
the charged defect's reverse-source contribution and the neutral
rank-filter fluctuation with the same original $J$, then resum
their mixed marked words in both memory times. The second-order
cross above is a concrete check on that construction. No bulk
mass assignment or finite-$G$ interaction law follows yet.
If $\lambda_N$ merely stays in a positive compact interval,
\eqref{eq:reduction} still holds by taking convergent subsequences;
explicit $q,C$ spectral formulas assume a chosen limit $\lambda$.

\begin{thebibliography}{9}\small
\bibitem{BKS} M. Bo\.{z}ejko, B. K\"ummerer and R. Speicher,
\emph{$q$-Gaussian processes: non-commutative and classical aspects},
\href{https://arxiv.org/abs/funct-an/9604010}{arXiv:funct-an/9604010},
Prop.~2.9, Theorems 2.11 and 2.13.
\bibitem{BNR} M. Berkooz, V. Narovlansky and H. Raj,
\emph{Complex Sachdev-Ye-Kitaev model in the double scaling limit},
\href{https://arxiv.org/abs/2006.13983}{arXiv:2006.13983}, Sec.~3.3.
\bibitem{Tall} J. Tall and S. Tomsovic,
\emph{Double-scaled bosonic and fermionic embedded ensembles,
complex SYK, and the dual Hilbert space},
\href{https://arxiv.org/abs/2604.14522}{arXiv:2604.14522v2}, Sec.~2.1.
\bibitem{Uniform} A. Day and Chengheng,
\emph{Uniform second moments of finite-time memory fields}, local note (7 Nov.~2026).
\bibitem{Return} A. Day and Chengheng,
\emph{The conserved background of the first memory return}, local note (6 Nov.~2026).
\bibitem{Central} A. Day and Chengheng,
\emph{Central charge windows in the memory spectrum}, local note (6 Nov.~2026).
\bibitem{Raw} A. Day and Chengheng,
\emph{The raw memory coefficient in a central charge window}, local stage
\texttt{memory-central-raw} (6 Nov.~2026).
\bibitem{Mixed} A. Day and Chengheng,
\emph{Mixed spectra of the first memory fields}, local note (6 Nov.~2026).
\end{thebibliography}
\end{document}
